\section{固定平面六次曲線と actual な個別 upper disks}
\label{ad:section}

本節は $t=\tau\ne0$ を固定する。
定理\ref{ai:c56-allends}の punctured curve を、元の pencil に存在する一つの閉曲線から選ぶ。
これにより自然な upper meridians を actual symplectic disks で個別に埋め、
その primitive の不足分を閉曲線の面積として計算できる。
Disk の族、cap fibration、あるいは既存積分の disk 全体への延長は主張しない。

\subsection{Punctured lift の具体的な選択と normal frames}

\begin{proposition}[固定曲線による actual lift]\label{ad:fixed-curve}
元の pencil
\[
 X_t=\left\{L_1L_2L_3L_4L_5L_6+tF_6=0\right\}\subset\PP^4
\]
に対し
\[
 C^*=\{L_5=L_6=F_6=0\}
\]
は全ての smooth $t\ne0$ の $X_t$ に含まれる同じ smooth plane sextic である。
その normal bundle と面積は
\begin{equation}\label{ad:normal-area}
 N_{C^*/X_t}\simeq\mathcal O_C(1)\oplus\mathcal O_C(1),
 \qquad \int_{C^*}H=6,\qquad \int_{C^*}\omega_t=6a
 \quad([\omega_t]=aH)
\end{equation}
である。$G=L_1L_2L_3L_4$ の零点24個を含む互いに素な小円板を削除すれば、
この punctured curve を既存 $\bar C_{56}$ の actual lift として選べる。
この選択は既存の明示的な collar family と一致する。
\end{proposition}

\begin{proof}
$L_5=L_6=0$ 上では pencil の方程式は $tF_6=0$ となる。
従って $t\ne0$ では $C^*$ は $X_t$ 内の二つの hyperplane sections の
scheme-theoretic な横断的 complete intersection である。
横断性は $F_6|_{P_5\cap P_6}$ の零曲線の smoothness から従い、
conormal sequence が normal bundle の式を与える。
Plane sextic の hyperplane degree は6なので面積の式も従う。
本稿の chosen class $aH+\eps A$ は $t\ne0$ 上では $aH$ に制限される。

各 puncture では、単元を吸収した既存の chart
\[
 xyz=tw,\qquad t=xT,\quad y=TY,\quad w=Yz
\]
を使う。$C^*_t$ の strict lift は正確に
\[
 Y=z=0,\qquad T=t/x
\]
である。これは式\eqref{sc:nonzero-chart}で用いた actual collar であり、
削除円板の外では $x\ne0$ のため $t=0$ へ smooth に延びる。
従って一般の horizontal lift を選ぶ代わりに、この固定曲線を選べばよい。
\end{proof}

Normal frames の次数差もこの同定で直接見える。
以下では順序を、$z=L_6$ 方向の $\mathcal O_C(1)$ と、
central $V_5$ の normal である $\mathcal O_C(5)$ の順とする。
元の pencil の係数で $\mathcal Y=F_6/L_6$ と書くと、$G\ne0$ 上では
\begin{equation}\label{ad:frame-comparison}
 y=L_5=-\frac{t}{G}\mathcal Y,
 \qquad
 \partial_{\mathcal Y}\longmapsto-\frac{t}{G}\partial_y
\end{equation}
が normal first jet の同定である。
ここで $\mathcal Y$ は central blow-up の $\mathcal O(5)$ 座標であり、
annular tube の $Y$ という名前と自動的には同一視しない。
式は $Gyz+tF_6=0$ を用いて $F_6/z$ を holomorphically 延長すれば得られる。
第一の $z$ 方向は変化しない。

$G|_C$ は24個の simple zeros を持つので、第二の frame comparison は
各 deleted disk の counterclockwise 境界で winding $-1$ を持つ。
従って normal degree は $(6,30)$ から $(6,6)$ へ変わる。
特に式\eqref{ai:individual-sections}の algebraic sections は
\begin{equation}\label{ad:actual-sections}
 (L_1,L_1L_2^2L_3L_4)\longmapsto(L_1,-tL_2)
\end{equation}
となる。Unitary normalization と orthogonalization はこの relative degree を変えない。
Semistable frame における determinant profile $(2,2,1,1)$ は、
actual disk frames では $(1,1,0,0)$ となる（各型6端）。
これは frame の比較であり、dual mixed Hessian の winding と向きを混同しない。

\subsection{個別の graph disks と exact な面積収支}

Actual $C^*\subset X_\tau$ 上で互いに素な閉円板 $\Delta_i$ を固定し、
その正向き境界を $c_i$、補集合を
\[
 C'=C^*\setminus\bigcup_{i=1}^{24}\operatorname{int}\Delta_i
\]
とする。各円板は対応する puncture を含み、$\partial C'=-\sum_i c_i$ である。
先に円板を十分小さく選び、
\[
 0<A_i^{\rm base}:=\int_{\Delta_i}\omega_\tau<a
\]
としてよい。以後は $\tau$、円板、collars、normal embeddings を固定する。
一つの global primitive $d\beta=\omega_\tau|_{C'}$ と既存の normal-model primitive を使い、
collars では
\[
 \beta=\widehat I_i\,d\theta_i,\qquad c_i=\{\widehat I_i=I_i^0\}
\]
とする。Actual primitive と model primitive の exact difference は periods に影響しない。

\begin{theorem}[Actual individual upper disks と次数6の budget]\label{ad:disks}
各 collar で固定した prototype parameters
\[
 H_0=h_i^0>0,\qquad J_0=0,\qquad
 c^2<q_{*,i}<(F_1(h_i^0,0)-c)^2
\]
と自然な upper meridian $\gamma_{U,i}$ を選ぶ。
Annular insertion の scale を $s_i>0$、offset を
\[
 b_i=I_i^0-s_i h_i^0
\]
とする。十分小さい scales に対し、各 actual $\gamma_{U,i}$ は
$\Delta_i$ 上の $C^1$-small normal graph である embedded symplectic disk $D_i$ を bound する。
その面積 $A_i$ と既存 primitive period $P_i$ は
\begin{align}
 P_i&=b_i+s_i B_i,\qquad B_i=B(h_i^0,0,q_{*,i}),\label{ad:period}\\
 A_i&=A_i^{\rm base}+P_i-I_i^0
     =A_i^{\rm base}-s_i(h_i^0-B_i),\label{ad:disk-area}\\
 0&<A_i<a.\nonumber
\end{align}
また、$c_i$ から $\gamma_{U,i}$ への radial normal annulus $\mathcal A_i$ を
$\partial\mathcal A_i=\gamma_{U,i}-c_i$ と向けると、closed cycle
\begin{equation}\label{ad:cycle}
 Z=C'+\sum_i(D_i-\mathcal A_i)
\end{equation}
は $[Z]=[C^*]\in H_2(X_\tau;\ZZ)$ を満たす。従って
\begin{equation}\label{ad:budget}
 \int_ZH=6,\qquad \boxed{\ \sum_{i=1}^{24}(A_i-P_i)=6a\ }.
\end{equation}
\end{theorem}

\begin{proof}
Twists を $(a_i,d_i)$ に配分した insertion を使う。
式\eqref{ai:distributed-end}の共通 curve-model coordinates では厳密に
\begin{equation}\label{ad:graph-coordinates}
 \widehat I_i=b_i+s_iH_0,\qquad
 U=\sqrt{s_i}\,e^{-2\pi i\theta_x}y,\qquad
 V=\sqrt{s_i}\,z,
 \qquad \theta_i=\theta_x.
\end{equation}
従って、ここでいう base projection は physical tube の $I$ ではなく
共通 global normal model の $\widehat I_i$ である。

自然な upper meridian では $\theta_y+\theta_z$ と $-\theta_z$ を固定するので、
$\theta_y,\theta_z$ は共に一定である。
$W=c+r e^{-2\pi i\ell}$, $r=\sqrt{q_{*,i}}>c$, と parameterize すれば
\[
 \frac{d}{d\ell}\arg W
 =-2\pi\frac{r(r+c\cos(2\pi\ell))}{|W|^2}<0.
\]
他方 $\theta_x=\theta_y+\theta_z-\arg W/(2\pi)$ だから、
$\theta_x$ は正向きに単調に一周する。
式\eqref{ad:graph-coordinates}では $\widehat I_i=I_i^0$ であり、
従って $\gamma_{U,i}$ は $c_i$ 上の normal section の graph である。
Prototype parameters と $\tau$ を固定しているため、その $C^1$ norm は $O(\sqrt{s_i})$ である。

閉曲線 $C^*$ の smooth tubular neighborhood を、選んだ compact collars 上では
既存の common normal embedding と一致するように選ぶ。
これは actual normal bundle を用いた relative tubular construction で可能である。
具体的には、二つの tubular germs の first jet を actual normal bundle で一致させ、
その差 $O(|v|^2)$ を base の cutoff で延ばせば、十分小さい radius で embedding のままである。
Disk 上では normal bundle は trivial なので、境界 section を
固定した collar extension と cutoff で $\Delta_i$ 全体へ延ばせる。
この extension の $C^1$ norm も $O(\sqrt{s_i})$ であり、中心近傍で零にしてよい。
得られる graph は positive complex disk $\Delta_i$ に $C^1$-close なので、
十分小さい $s_i$ では embedded symplectic disk となる。
Frame 自身が disk へ nonvanishing に延びる必要はなく、section は零点を持ってよい。
有限個の円板について同時に scales を小さくできる。

Radial normal annulus は既存の primitive が定義された collar tube 内にあり、Stokes により
\[
 \int_{\mathcal A_i}\omega_\tau=P_i-I_i^0.
\]
Disk の normal section を零へ縮める homotopy は
$D_i-\Delta_i-\mathcal A_i$ を境界とする3-chain を与える。
従って式\eqref{ad:disk-area}と $[Z]=[C^*]$ が従う。
補題\ref{cf:cap-area-bound}の $0<B_i<h_i^0/3$ より、
$s_i(h_i^0-B_i)<A_i^{\rm base}$ まで scale を縮めれば
$0<A_i<A_i^{\rm base}<a$ である。

最後に $\operatorname{Area}(C')=-\sum_iI_i^0$ なので
\[
 \int_Z\omega_\tau
 =\operatorname{Area}(C')+
   \sum_i\bigl(A_i-P_i+I_i^0\bigr)
 =\sum_i(A_i-P_i)=6a.
\]
これで向きと primitive constants を含む budget が得られる。
\end{proof}

\paragraph{範囲と central limit.}
この構成は各指定 cycle を埋める個別の disk を与える。
Disk 内部と既存 assembly 全体との非交差、$T^2$-family of disks、
Lagrangian cap fibration、所定の $H,J,Q$ の延長はこの定理の結論に含めない。
また既存 primitive が $D_i$ 全体へ延びるとも主張しない。
その差の総和が $6a$ であるため、式\eqref{cf:negative-cap-period}との矛盾はない。
小ささの定数は固定した $\tau\ne0$ に依存し、$\tau\to0$ の一様な radius は得ていない。

閉曲線族の central limit も、punctured lift の limit と区別する必要がある。
各 puncture の chart では $Y=z=0$, $t=xT$ なので、$t=0$ には
strict $C_{56}$ に対応する $T=0$ と、最初の blow-up の exceptional curve に対応する $x=0$ が現れる。
従って closed limit cycle は $C_{56}$ 単独ではなく、
24個の exceptional $\PP^1$ を付加したものである。
式\eqref{ad:normal-area}の閉曲線の面積を、strict $C_{56}$ の面積だけと同一視してはいけない。
